\documentclass[11pt,a4paper]{article}
\usepackage[margin=2.5cm]{geometry}
\usepackage[T1]{fontenc}
\usepackage{booktabs}
\usepackage{xcolor}
\usepackage{amsmath}
\usepackage{listings}
\usepackage[hidelinks]{hyperref}

\emergencystretch=2em
\definecolor{codebg}{gray}{0.95}
\lstset{language=C++, basicstyle=\ttfamily\small, keywordstyle=\color{blue!70!black},
  commentstyle=\color{green!50!black}\itshape, backgroundcolor=\color{codebg},
  frame=single, rulecolor=\color{black!20}, breaklines=true, showstringspaces=false,
  tabsize=2, columns=flexible}

\title{ATmega32A Module \\ \Large Timers (CTC Interrupt Mode)}
\author{Ali Sahafi}
\date{}

\begin{document}
\maketitle

\section{Theory}

A hardware timer is a counter that increments automatically with every clock
tick --- independently of your program. The ATmega32A has three:

\begin{center}
\begin{tabular}{llll}
\toprule
\textbf{Timer} & \textbf{Width} & \textbf{Max count} & \textbf{Compare register} \\
\midrule
Timer0 & 8-bit  & 255    & \texttt{OCR0}  \\
Timer1 & 16-bit & 65535  & \texttt{OCR1A} \\
Timer2 & 8-bit  & 255    & \texttt{OCR2}  \\
\bottomrule
\end{tabular}
\end{center}

In \textbf{CTC mode} (\emph{Clear Timer on Compare match}) the counter runs
from 0 up to the value in the compare register, fires an interrupt, resets to
0, and starts again. Because the CPU clock is usually too fast to count
directly, a \textbf{prescaler} divides it first. The interrupt frequency is:

\begin{equation*}
f_{\text{interrupt}} = \frac{F_{\text{CPU}}}{\text{prescaler} \cdot (\text{OCR} + 1)}
\end{equation*}

\paragraph{Worked example.} A 1-second tick at $F_{\text{CPU}} = 8$\,MHz:
choose prescaler 256, then $\text{OCR} = \frac{8\,000\,000}{256 \cdot 1} - 1
= 31249$. That exceeds 255, so it only fits the 16-bit Timer1 --- exactly
what the example below uses.

\section{Registers}

\begin{center}
\small
\begin{tabular}{lll}
\toprule
\textbf{Register} & \textbf{Role} & \textbf{Bits used by the driver} \\
\midrule
\texttt{TCCR0}, \texttt{TCCR2}   & Timer0/2 control & \texttt{WGM$x$1} (CTC), \texttt{CS} bits (prescaler) \\
\texttt{TCCR1A}, \texttt{TCCR1B} & Timer1 control   & \texttt{WGM12} (CTC), \texttt{CS} bits (prescaler) \\
\texttt{OCR0}, \texttt{OCR1A}, \texttt{OCR2} & Compare value & the ``count-to'' number \\
\texttt{TCNT0}, \texttt{TCNT1}, \texttt{TCNT2} & Counter value & readable/resettable at any time \\
\texttt{TIMSK} & Interrupt mask & \texttt{OCIE0}, \texttt{OCIE1A}, \texttt{OCIE2} \\
\bottomrule
\end{tabular}
\end{center}

\section{API reference}

\begin{center}
\small
\begin{tabular}{ll}
\toprule
\textbf{Method} & \textbf{Description} \\
\midrule
\texttt{Timer.initTimer0(ocr, prescaler)} & 8-bit CTC, OCR max 255 \\
\texttt{Timer.initTimer1(ocr, prescaler)} & 16-bit CTC, OCR max 65535 \\
\texttt{Timer.initTimer2(ocr, prescaler)} & 8-bit CTC, OCR max 255 \\
\texttt{Timer.attachTimer0/1/2(callback)} & Set callback: \texttt{void cb()} \\
\texttt{Timer.stop0/1/2()} & Halt counter (preserves config) \\
\texttt{Timer.start0/1/2()} & Resume from current count \\
\texttt{Timer.reset0/1/2()} & Set counter to 0 (keeps running) \\
\texttt{Timer.restart0/1/2()} & Set counter to 0 and start \\
\texttt{Timer.read0/2()} & Read 8-bit counter (0--255) \\
\texttt{Timer.read1()} & Read 16-bit counter (0--65535) \\
\texttt{Timer.delay\_ms(ms)} & Blocking delay (no interrupt needed) \\
\bottomrule
\end{tabular}
\end{center}

\textbf{Prescalers:} \texttt{TIMER0\_PRESCALER\_1/8/64/256/1024},
\texttt{TIMER1\_PRESCALER\_1/8/64/256/1024},
\texttt{TIMER2\_PRESCALER\_1/8/32/64/128/256/1024}.

\paragraph{Warning: Timer2 encoding.} Timer2's prescaler \emph{bit encoding}
differs from Timer0/1 (it offers 32 and 128 as extra options). Always use the
matching \texttt{TIMER2\_PRESCALER\_xxx} constants for Timer2 --- e.g.
\texttt{TIMER0\_PRESCALER\_64} and \texttt{TIMER2\_PRESCALER\_64} are
\emph{different} register values.

\paragraph{Warning: PWM conflict.} Timer0 and Timer1 are shared with the PWM
module. Never call \texttt{Timer.initTimer0/1()} and
\texttt{PWM.initTimer0/1\_FastPWM()} in the same program --- they program the
same \texttt{TCCR} registers.

\paragraph{Don't forget \texttt{sei()}.} Timer interrupts only fire after
global interrupts are enabled.

\section{Example: 1-second LED blink, zero CPU load}

\paragraph{Wiring.}
\begin{itemize}
  \item LED + 330\,$\Omega$ resistor from \textbf{PC0} to GND.
\end{itemize}

\paragraph{Build \& flash.} \texttt{make timer}

\lstinputlisting{example.cpp}

The main loop is empty --- the blink happens entirely ``in the background''
via the interrupt. Callbacks run in interrupt context: keep them short, and
declare variables shared with \texttt{main()} as \texttt{volatile}.

\section{Exercises}

\begin{enumerate}
  \item Make the LED blink every 250\,ms instead of every second. Compute the
        new OCR value (keep prescaler 256) before touching the code.
  \item Compute OCR and prescaler for a \textbf{1\,ms} tick on
        \textbf{Timer0} at 8\,MHz. Then build a stopwatch: count milliseconds
        in the callback and print seconds on the seven-segment display.
  \item Run Timer1 (1\,s, LED on PC0) and Timer2 ($\approx$31\,ms, LED on PC1)
        simultaneously and observe that the two blink patterns are
        independent.
\end{enumerate}

\end{document}
